\begin{proposition}[H\"older's inequality on a set of finite measure]
Let $(X,\mu)$ be a measure space, let $E\subseteq X$ be measurable, and let
$q,q'\in(1,\infty)$ be conjugate, $\tfrac1q+\tfrac1{q'}=1$.  If $f,g\ge0$ are
measurable on $E$ with $f\in L^{q}(E)$ and $g\in L^{q'}(E)$, then
\[
\int_E fg\,d\mu\ \le\ \Bigl(\int_E f^{q}\,d\mu\Bigr)^{1/q}
\Bigl(\int_E g^{q'}\,d\mu\Bigr)^{1/q'} .
\]
Consequently, if $\mu(E)<\infty$ and $q\in[1,\infty)$, then taking $f\equiv1$
gives, for every measurable $g$ with $|g|\in L^{q}(E)$,
\[
\int_E|g|\,d\mu\ \le\ \mu(E)^{1-\frac1q}\Bigl(\int_E|g|^{q}\,d\mu\Bigr)^{1/q}
\ <\ \infty ,
\]
so $L^{q}(E)\subseteq L^{1}(E)$ whenever $\mu(E)<\infty$ and $q\ge1$.  More
generally $L^{q}(E)\subseteq L^{s}(E)$ for $1\le s\le q$ and $\mu(E)<\infty$.
In Mathlib the inequality is
\texttt{MeasureTheory.integral\_mul\_le\_Lp\_mul\_Lq\_of\_nonneg}; the
quantitative inclusion with its $\mu(E)^{1/s-1/q}$ factor is
\texttt{MeasureTheory.eLpNorm\_le\_eLpNorm\_mul\_rpow\_measure\_univ}, and the
membership statement $L^{q}\subseteq L^{s}$ for a finite measure is
\texttt{MeasureTheory.MemLp.mono\_exponent}.  The last two are applied to the
restricted measure $\mu|_E$, which is finite exactly when $\mu(E)<\infty$.

Two further facts about $L^{q}$ of a set, used wherever a Sobolev-type
membership is transported between sets, are recorded here rather than repeated.
First, $L^{q}$ membership is inherited by subsets: if $E'\subseteq E$ are
measurable and $f\in L^{q}(E)$ then $f\in L^{q}(E')$, because
$\mu|_{E'}\le\mu|_{E}$ and $L^{q}$ membership is monotone in the measure
(\texttt{MeasureTheory.Measure.restrict\_mono} and
\texttt{MeasureTheory.MemLp.mono\_measure}).  Second, $L^{q}$ of a set is the
same thing as $L^{q}$ of the zero-extension: for measurable $E$,
\[
f\in L^{q}(E)\iff \mathbf1_Ef\in L^{q}(X)
\]
(\texttt{MeasureTheory.memLp\_indicator\_iff\_restrict}, whose only hypothesis
is measurability of $E$; the underlying norm identity is
\texttt{MeasureTheory.eLpNorm\_indicator\_eq\_eLpNorm\_restrict}).  Combining the
two: if $f$ vanishes outside a measurable $E\subseteq E''$ and $f\in L^{q}(E)$,
then $f\in L^{q}(E'')$, since $\mathbf1_Ef=f$ and
$\mathbf1_{E''}f$ have the same $L^q(X)$ norm.
\end{proposition}

\begin{proposition}[The $L^{q}$ triangle inequality (Minkowski)]
Let $(X,\mu)$ be a measure space and $q\in[1,\infty]$.  If $f,g$ are measurable
then
\[
\|f+g\|_{L^{q}(\mu)}\ \le\ \|f\|_{L^{q}(\mu)}+\|g\|_{L^{q}(\mu)} ,
\]
and in particular $L^{q}(\mu)$ is closed under addition: $f,g\in L^{q}(\mu)$
implies $f+g\in L^{q}(\mu)$.  It is also closed under multiplication by a real
scalar, $\|cf\|_{L^q(\mu)}=|c|\,\|f\|_{L^q(\mu)}$, so $f\in L^q(\mu)$ implies
$cf\in L^q(\mu)$; consequently $L^q(\mu)$ is closed under real linear
combinations.  In Mathlib these are
\texttt{MeasureTheory.eLpNorm\_add\_le} (which needs $1\le q$ and almost
everywhere strong measurability of $f$ and $g$),
\texttt{MeasureTheory.MemLp.add}, and
\texttt{MeasureTheory.MemLp.const\_mul} / \texttt{MeasureTheory.MemLp.const\_smul}
for the scalar multiples.  All of them hold for functions valued in a real
Banach space as well as for real-valued ones, which is how they are applied to
weak gradients below.  A constant function lies in $L^q(\mu)$ when $\mu$ is
finite (\texttt{MeasureTheory.memLp\_const}).
\end{proposition}

\begin{proposition}[Cauchy--Schwarz]
Let $F$ be a real inner product space.  Then for all $x,y\in F$,
\[
|(x,y)|\ \le\ |x|\,|y| ,
\qquad\text{equivalently}\qquad
(x,y)^2\ \le\ (x,x)\,(y,y) .
\]
In Mathlib these are \texttt{abs\_real\_inner\_le\_norm} and
\texttt{real\_inner\_mul\_inner\_self\_le}.  Both apply to
$\R^n=\texttt{EuclideanSpace}\ \R\ (\texttt{Fin}\ n)$ with its standard inner
product, for which $(x,y)=\sum_{i}x_iy_i$ and $|x|^2=(x,x)$.
\end{proposition}

\begin{proposition}[A nonnegative function with vanishing integral vanishes almost everywhere]
Let $(X,\mu)$ be a measure space and let $f:X\to\R$ be integrable with
$f\ge0$ $\mu$-almost everywhere.  Then
\[
\int_X f\,d\mu=0\quad\Longleftrightarrow\quad f=0\ \ \mu\text{-a.e.}
\]
In Mathlib this is
\texttt{MeasureTheory.integral\_eq\_zero\_iff\_of\_nonneg\_ae}.  Note that both
hypotheses are needed: integrability of $f$ and almost everywhere
nonnegativity.  Applied on a measurable subset $E$ with the restricted measure,
it says that an integrable $f\ge0$ with $\int_Ef\,d\mu=0$ vanishes a.e.\ on
$E$.
\end{proposition}

\begin{proposition}[Lebesgue dominated convergence]
Let $(X,\mu)$ be a measure space, let $G$ be a real Banach space, and let
$F_k:X\to G$ be almost everywhere strongly measurable for every $k\in\mathbb{N}$.
Suppose there is an integrable $b:X\to\R$ with $|F_k|\le b$ $\mu$-a.e.\ for
every $k$, and suppose $F_k\to f$ $\mu$-a.e.  Then
\[
\lim_{k\to\infty}\int_X F_k\,d\mu=\int_X f\,d\mu .
\]
In Mathlib this is
\texttt{MeasureTheory.tendsto\_integral\_of\_dominated\_convergence}, whose
four hypotheses are exactly measurability of each $F_k$, integrability of the
bound $b$, the almost everywhere domination $|F_k|\le b$ for each $k$, and the
almost everywhere convergence $F_k\to f$.
\end{proposition}

\begin{proposition}[Measurability of superlevel sets and of countable operations on them]
Let $X$ be a measurable space and let $f:X\to\R$ be measurable.  Then for every
$t\in\R$ the superlevel set
\[
\{f>t\}=\{x\in X: t<f(x)\}
\]
is measurable; the same holds for $\{f<t\}$, and hence for $\{f\ge t\}$ and
$\{f\le t\}$ by complementation.  Moreover the measurable sets are closed under
differences $S_1\setminus S_2$ and under countable unions and intersections
$\bigcup_{k}S_k$, $\bigcap_{k}S_k$.  In Mathlib the superlevel set is
\texttt{measurableSet\_lt} (applied with the constant function $t$ and with
$f$), and the set operations are \texttt{MeasurableSet.diff},
\texttt{MeasurableSet.iUnion} and \texttt{MeasurableSet.iInter}.
\end{proposition}

\begin{proposition}[The fundamental theorem of calculus in interval-integral form]
Let $E$ be a real Banach space, let $a,b\in\R$, and let
$\varphi,\varphi':\R\to E$ be such that $\varphi$ has derivative $\varphi'(t)$
at every point $t$ of the closed interval with endpoints $a$ and $b$.  If
$\varphi'$ is interval-integrable on that interval, then
\[
\int_a^b\varphi'(t)\,dt=\varphi(b)-\varphi(a) .
\]
In Mathlib this is
\texttt{intervalIntegral.integral\_eq\_sub\_of\_hasDerivAt}.  There is also a
variant requiring differentiability only on the open interval, with continuity
on the closed interval,
\texttt{intervalIntegral.integral\_eq\_sub\_of\_hasDerivAt\_of\_le}.  Applied to
$\varphi(t)=\Phi(tx+(1-t)y)$ for $\Phi$ differentiable on $\R^n$, this is the
segment form used repeatedly below: $\Phi(x)-\Phi(y)=\int_0^1
D\Phi\bigl(tx+(1-t)y\bigr)(x-y)\,dt$, valid as soon as the integrand is
interval-integrable on $[0,1]$.
\end{proposition}

\begin{proposition}[Interval integrability of $t^{r}$ for $r>-1$]
Let $r\in\R$ with $r>-1$ and let $a,b\in\R$.  Then $t\mapsto t^{r}$, the real
power \texttt{Real.rpow}, is interval-integrable with respect to Lebesgue
measure on the interval with endpoints $a$ and $b$; in particular $\int_0^1t^{r}\,dt=\frac1{r+1}<\infty$.
In Mathlib this is \texttt{intervalIntegral.intervalIntegrable\_rpow'}, whose
only hypothesis is $-1<r$, so the interval may straddle or lie below $0$.  The
conventions of \texttt{Real.rpow} off the positive axis are $0^{r}=0$ for
$r\ne0$ and $0^{0}=1$ (\texttt{Real.zero\_rpow}), and
$t^{r}=|t|^{r}\cos(r\pi)$ for $t<0$ (\texttt{Real.rpow\_def\_of\_neg}) -- in
particular $t^{r}$ is \emph{not} $0$ for $t<0$ in general.  The hypothesis
$r>-1$ is sharp: for $r\le-1$ the integral over $[0,1]$ diverges.
\end{proposition}

\begin{proposition}[Measurability of the derivative]
Let $E,F$ be finite-dimensional real normed spaces and let $\Phi:E\to F$ be an
arbitrary map.  Then $x\mapsto D\Phi(x)$ is measurable, where $D\Phi(x)$ is
taken to be $0$ at points of non-differentiability; consequently, for each fixed
$y\in E$ the map $x\mapsto D\Phi(x)y$ is measurable, and for $F$-valued
$\psi:\R\to F$ the map $t\mapsto\psi'(t)$ is strongly measurable.  No
differentiability or continuity of $\Phi$ is assumed.  In Mathlib these are
\texttt{measurable\_fderiv}, \texttt{measurable\_fderiv\_apply\_const} and
\texttt{stronglyMeasurable\_deriv}.  This is what makes a derivative obtained
from bare differentiability -- as opposed to $C^1$ regularity -- eligible for
the integrability hypothesis of the fundamental theorem of calculus above.
\end{proposition}

\begin{proposition}[Domination, monotonicity and the norm inequality for the integral]
Let $(X,\mu)$ be a measure space and let $G$ be a real Banach space.
\emph{(a) Domination.}  If $g:X\to\R$ is integrable, $f:X\to G$ is almost
everywhere strongly measurable and $|f|\le g$ $\mu$-a.e., then $f$ is
integrable.  \emph{(b) Monotonicity.}  If $f,g:\R\to\R$ are
interval-integrable on $[\alpha,\beta]$ with $\alpha\le\beta$ and $f\le g$
there, then $\int_\alpha^\beta f\le\int_\alpha^\beta g$.  \emph{(c) Norm
inequality.}  If $f:\R\to G$ is interval-integrable on $[\alpha,\beta]$ with
$\alpha\le\beta$, then
\[
\Bigl|\int_\alpha^\beta f(t)\,dt\Bigr|\ \le\ \int_\alpha^\beta|f(t)|\,dt .
\]
In Mathlib these are \texttt{MeasureTheory.Integrable.mono'} together with its
interval form \texttt{IntervalIntegrable.mono\_fun},
\texttt{intervalIntegral.integral\_mono\_on} (and its almost everywhere variant
\texttt{intervalIntegral.integral\_mono\_ae}), and
\texttt{intervalIntegral.norm\_integral\_le\_integral\_norm}.
\end{proposition}

\begin{proposition}[Density of the compactly supported continuous functions in $L^{q}$]
Let $\mu$ be a (weakly regular, $\sigma$-finite) Borel measure on $\R^n$ --- in
particular Lebesgue measure --- let $q\in[1,\infty)$ and let $w\in L^{q}(\mu)$.
Then for every $\sigma>0$ there is a continuous $v\colon\R^n\to\R$ with compact
support, $v\in L^{q}(\mu)$, such that
\[
\|w-v\|_{L^{q}(\mu)}\ \le\ \sigma .
\]
In Mathlib this is
\texttt{MeasureTheory.MemLp.exists\_hasCompactSupport\_eLpNorm\_sub\_le}, which
returns a $v$ that is simultaneously continuous, compactly supported and in
$L^{q}$.
\end{proposition}

\begin{proposition}[Invariance of Lebesgue measure under translations and unimodular linear maps]
Lebesgue measure on $\R^n$ is invariant under every translation
$x\mapsto x+a$: for measurable $S$ one has $|S+a|=|S|$, and consequently, for
every $q\in[1,\infty]$ and every almost everywhere strongly measurable
$w\colon\R^n\to\R$,
\[
\|w(\cdot+a)\|_{L^{q}(\R^n)}=\|w\|_{L^{q}(\R^n)} ,
\qquad
\int_{S+a}w(x)\,dx=\int_{S}w(x+a)\,dx .
\]
In Mathlib the measure statement is
\texttt{MeasureTheory.measure\_preimage\_add} together with
\texttt{measurePreserving\_add\_right} (Lebesgue measure is both left and right
invariant, being a Haar measure of an abelian group), and the seminorm statement
follows by \texttt{MeasureTheory.eLpNorm\_comp\_measurePreserving} applied to
that measure-preserving map.

More generally, if $T\colon\R^N\to\R^N$ is linear with $\det T\ne0$ then
$|T^{-1}(S)|=|\det T|^{-1}|S|$ for every measurable $S$, and the pushforward of
Lebesgue measure under $T$ is $|\det T|^{-1}$ times Lebesgue measure; in
particular, when $|\det T|=1$ the map $T$ preserves Lebesgue measure, hence
preserves Lebesgue measurability and null sets, and
$\int_{\R^N}w\circ T=\int_{\R^N}w$ for every nonnegative measurable $w$.  In
Mathlib these are
\texttt{MeasureTheory.Measure.addHaar\_preimage\_linearMap} and
\texttt{MeasureTheory.Measure.map\_linearMap\_addHaar\_eq\_smul\_addHaar}.  The
case used below is the shear $T(x,s)=(x+s\,e_k,\ s)$ of
$\R^n\times\R=\R^{n+1}$, whose matrix is triangular with unit diagonal, so
$\det T=1$.
\end{proposition}

\begin{proposition}[Tonelli and Fubini]
Let $\mu,\nu$ be $\sigma$-finite measures and let $\Theta\ge0$ be measurable on
the product.  Then the two iterated integrals agree in $[0,\infty]$:
\[
\int\!\!\Bigl(\int\Theta(x,y)\,d\nu(y)\Bigr)d\mu(x)
=\int\!\!\Bigl(\int\Theta(x,y)\,d\mu(x)\Bigr)d\nu(y) .
\]
No integrability hypothesis is needed, only measurability on the product and
nonnegativity; both sides may be $+\infty$.  In Mathlib this is
\texttt{MeasureTheory.lintegral\_lintegral\_swap}.

The same exchange holds for a Banach-space valued $\Theta$ that is integrable on
the product rather than nonnegative, which is Fubini's theorem,
\texttt{MeasureTheory.integral\_integral\_swap}.  A continuous function with
compact support on $\R^N$ is integrable, so the exchange is available for such
$\Theta$.
\end{proposition}

\begin{proposition}[An almost everywhere convergent subsequence out of $L^{q}$ convergence]
Let $q\in[1,\infty)$, let $w_j,w$ be almost everywhere strongly measurable with
$\|w_j-w\|_{L^{q}(\mu)}\to0$.  Then there is a strictly increasing sequence
$j_1<j_2<\cdots$ with $w_{j_i}\to w$ $\mu$-almost everywhere.  In Mathlib this
is \texttt{MeasureTheory.tendstoInMeasure\_of\_tendsto\_eLpNorm\_of\_ne\_top}
(which turns the $L^{q}$ convergence into convergence in measure, for
$0\ne q\ne\infty$) followed by
\texttt{MeasureTheory.TendstoInMeasure.exists\_seq\_tendsto\_ae}.
\end{proposition}

\begin{proposition}[Convolution with a smooth compactly supported kernel]
Let $w\in L^{1}_{\mathrm{loc}}(\R^n)$ and let $\rho\colon\R^n\to\R$ be of class
$C^{\infty}$ with compact support.  Then the convolution
\[
(w*\rho)(x)=\int_{\R^n}w(y)\,\rho(x-y)\,dy
\]
is of class $C^{\infty}$ on $\R^n$, and its partial derivatives are obtained by
differentiating the kernel:
\[
D_i(w*\rho)=w*(D_i\rho)\qquad(i=1,\dots,n).
\]
In Mathlib these are \texttt{HasCompactSupport.contDiff\_convolution\_right}
(smoothness, from local integrability of $w$ and smoothness and compact support
of $\rho$) and \texttt{HasCompactSupport.hasFDerivAt\_convolution\_right}, whose
conclusion is that $w*\rho$ is differentiable at every point with differential
the convolution of $w$ against the differential of $\rho$.
\end{proposition}

\begin{proposition}[The fundamental lemma of the calculus of variations (du Bois--Reymond)]
Let $U\subseteq\R^n$ be open, let $F$ be a real Banach space and let
$w\colon\R^n\to F$ be locally integrable on $U$.  Suppose that
\[
\int_{\R^n}\varphi(x)\,w(x)\,dx=0
\]
for every $\varphi\in C^\infty(\R^n)$ with compact support and
$\operatorname{supp}\varphi\subseteq U$.  Then $w=0$ almost everywhere on $U$.
In Mathlib this is
\texttt{IsOpen.ae\_eq\_zero\_of\_integral\_contDiff\_smul\_eq\_zero}, whose
hypotheses are exactly openness of $U$, \texttt{LocallyIntegrableOn} of $w$ on
$U$, and the vanishing of $\int\varphi\cdot w$ for every $\varphi$ that is
$C^\infty$, compactly supported and with $\operatorname{tsupport}$ inside $U$ ---
\emph{verbatim} the test class of \noderef{HasWeakGradientOn}; the conclusion is
the ``almost every $x$, if $x\in U$ then $w(x)=0$'' form.  The integral in the
hypothesis is over all of $\R^n$, which for an integrand vanishing off
$\operatorname{supp}\varphi\subseteq U$ is the same as the integral over $U$.
The version without the open set, \texttt{ae\_eq\_zero\_of\_integral\_contDiff\_smul\_eq\_zero},
asks for global local integrability and drops the support restriction.

Two consequences used below.  Applied to a difference $w_1-w_2$ of locally
integrable functions, it says that two such functions with the same integrals
against all test functions agree almost everywhere on $U$.  Applied to the
weak-gradient identity of \noderef{HasWeakGradientOn}, it says that a weak
gradient is determined almost everywhere on $U$, and in particular that a
function whose weak gradient is tested only against functions supported in an
open subset on which it vanishes has that weak gradient vanishing almost
everywhere there.
\end{proposition}

\begin{proposition}[Symmetry of the second derivative (Clairaut)]
Let $E,F$ be real normed spaces and let $\Phi\colon E\to F$ be differentiable at
every point, with derivative $D\Phi\colon E\to\mathcal L(E,F)$.  If $D\Phi$ is
itself differentiable at a point $x$, with derivative
$D^2\Phi(x)\in\mathcal L(E,\mathcal L(E,F))$, then $D^2\Phi(x)$ is symmetric:
\[
D^2\Phi(x)[v][w]=D^2\Phi(x)[w][v]\qquad\text{for all }v,w\in E .
\]
In Mathlib this is \texttt{second\_derivative\_symmetric}, with the local variant
\texttt{second\_derivative\_symmetric\_of\_eventually} requiring the first
derivative only on a neighbourhood of $x$.  Note that the hypothesis is
\emph{Fr\'echet twice differentiability} and nothing more --- no $C^2$
regularity is required.  For $\Phi\in C^\infty(\R^n)$ this gives the classical
$D_iD_j\Phi=D_jD_i\Phi$ at every point.
\end{proposition}

\begin{proposition}[Integration by parts in a direction]
Let $f,g\colon\R^n\to\R$ and let $v\in\R^n$.  Assume $f$ is differentiable at
every point of $\operatorname{supp}g$, that $g$ is differentiable at every point
of $\operatorname{supp}f$, and that the three functions
$x\mapsto D_vf(x)\,g(x)$, $x\mapsto f(x)\,D_vg(x)$ and $x\mapsto f(x)g(x)$ are
Lebesgue integrable on $\R^n$, where $D_v$ denotes the directional derivative
$D_vf(x)=Df(x)[v]$.  Then
\[
\int_{\R^n}f(x)\,D_vg(x)\,dx=-\int_{\R^n}D_vf(x)\,g(x)\,dx .
\]
In Mathlib this is
\texttt{integral\_mul\_fderiv\_eq\_neg\_fderiv\_mul\_of\_integrable}, stated for
a Haar measure on a finite-dimensional real normed space with the Borel
$\sigma$-algebra --- hypotheses that Lebesgue measure on
$\texttt{EuclideanSpace}\ \R\ (\texttt{Fin}\ n)$ satisfies --- and with exactly
the three integrability assumptions and the two differentiability assumptions
listed above.  The typical application here takes $v=e_i$ the $i$-th standard
basis vector, so that $D_v=D_i$, and one of $f,g$ continuous with compact
support, in which case all three integrability hypotheses follow because a
continuous function with compact support is integrable and the product of a
continuous function with a continuous compactly supported one is again
continuous with compact support.
\end{proposition}

\begin{proposition}[The volume of a ball]
Let $n\ge1$, let $x_0\in\R^n$ and let $\rho\ge0$.  Then the Lebesgue measure of
the ball $B_\rho(x_0)$ does not depend on its centre and scales as the $n$-th
power of its radius:
\[
|B_\rho(x_0)|=|B_\rho(0)|=\rho^{\,n}\,|B_1(0)| .
\]
Writing $\omega_n=|B_1(0)|\in(0,\infty)$, this is $|B_\rho(x_0)|=\omega_n\rho^n$;
in particular $|B_{2\rho}|=2^n|B_\rho|$ and
$|B_{2\rho}\setminus B_\rho|=(1-2^{-n})\,|B_{2\rho}|$.  In Mathlib these are
\texttt{MeasureTheory.Measure.addHaar\_ball\_center} (independence of the
centre) and \texttt{MeasureTheory.Measure.addHaar\_ball} (the factor
$\rho^{\operatorname{finrank}}$), both for a Haar measure on a
finite-dimensional real normed space, hypotheses that Lebesgue measure on
$\texttt{EuclideanSpace}\ \R\ (\texttt{Fin}\ n)$ satisfies; the second needs
$\rho\ge0$ and a nontrivial space.
\end{proposition}

\begin{proposition}[A standard mollifier]
Let $n\ge1$ and $\delta>0$.  There is a function $\rho_\delta\colon\R^n\to\R$
with
\[
\rho_\delta\in C^\infty(\R^n),
\qquad
\rho_\delta\ge0,
\qquad
\operatorname{supp}\rho_\delta\subseteq\overline{B_\delta(0)}\ \text{ compact},
\qquad
\int_{\R^n}\rho_\delta\,dz=1 .
\]
In Mathlib such a family is provided by \texttt{ContDiffBump}: for a bump $f$
centred at $0$ with outer radius $\delta$, the normalisation $f.\mathtt{normed}$
is smooth (\texttt{ContDiffBump.contDiff\_normed}), nonnegative
(\texttt{ContDiffBump.nonneg\_normed}), has integral $1$
(\texttt{ContDiffBump.integral\_normed}), has
$\operatorname{tsupport}$ exactly $\overline{B_\delta(0)}$
(\texttt{ContDiffBump.tsupport\_normed\_eq}) and compact support
(\texttt{ContDiffBump.hasCompactSupport\_normed}).  Two further properties of
the same object are used below, and they too need no formula.  Take the bump with
inner radius $\tfrac78\delta$ and outer radius $\delta$.  On the inner ball it
is identically $1$ (\texttt{ContDiffBump.one\_of\_mem\_closedBall}), so by
$\rho_\delta=f/\!\int f$ (\texttt{ContDiffBump.normed\_def}) and
$\int f>0$ (\texttt{ContDiffBump.integral\_pos}) it is there equal to a positive
constant:
\[
\rho_\delta(z)=\kappa_\delta>0
\qquad\text{for }|z|\le\tfrac78\delta ,
\]
and everywhere it is bounded by the reciprocal of the measure of that inner ball
(\texttt{ContDiffBump.normed\_le\_div\_measure\_closedBall\_rIn}):
\[
\sup_{\R^n}\rho_\delta\ \le\ \frac1{\bigl|B_{7\delta/8}(0)\bigr|}<\infty .
\]
These seven properties are the only ones used wherever a mollifier appears
below; no explicit formula for $\rho_\delta$ is needed, and none should be
relied on.
\end{proposition}

\begin{proposition}[Fatou's lemma]
Let $(X,\mu)$ be a measure space and let $w_m\colon X\to[0,\infty]$ be
measurable for every $m\in\mathbb{N}$.  Then
\[
\int_X\liminf_{m\to\infty}w_m\,d\mu\ \le\ \liminf_{m\to\infty}\int_Xw_m\,d\mu .
\]
In particular, if $w_m\ge0$ are measurable and $w_m\to w$ $\mu$-almost
everywhere, then $\int_Xw\,d\mu\le\liminf_m\int_Xw_m\,d\mu$: no uniform bound
and no integrability of the limit are assumed, and both sides may be $+\infty$.
In Mathlib this is \texttt{MeasureTheory.lintegral\_liminf\_le} for measurable
$w_m$ and \texttt{MeasureTheory.lintegral\_liminf\_le'} for almost everywhere
measurable ones, stated for the $[0,\infty]$-valued lower Lebesgue integral,
which for a nonnegative measurable function agrees with the Bochner integral
whenever the latter is defined.
\end{proposition}

\begin{proposition}[The essential supremum as the limit of normalised $L^q$ means]
Let $(X,\mu)$ be a measure space with $0<\mu(X)<\infty$, let
$f\colon X\to\R$ be almost everywhere strongly measurable and put
$M=\operatorname*{ess\,sup}_X|f|\in[0,\infty]$.  For $q\in[1,\infty)$ write
\[
N_q=\Bigl(\frac1{\mu(X)}\int_X|f|^{q}\,d\mu\Bigr)^{1/q}.
\]
Then:
\begin{enumerate}
\item $N_q\le M$ for every $q$.  This is monotonicity of the integral
  (Preamble) applied to the almost everywhere bound $|f|\le M$, which is
  \texttt{ae\_le\_essSup}.
\item If $0\le t<M$ then $\mu\{|f|>t\}>0$, and
  $N_q\ge t\bigl(\mu\{|f|>t\}/\mu(X)\bigr)^{1/q}$ for every $q$.  The first
  assertion is the contrapositive of \texttt{essSup\_le\_of\_ae\_le} --- if
  $\mu\{|f|>t\}=0$ then $|f|\le t$ almost everywhere and so $M\le t$ --- and the
  second is the Chebyshev--Markov inequality
  \texttt{MeasureTheory.pow\_mul\_meas\_ge\_le\_eLpNorm}, which bounds
  $\bigl(\varepsilon\,\mu\{\varepsilon\le|f|^{q}\}\bigr)^{1/q}$ below the
  $L^{q}$ seminorm, taken at $\varepsilon=t^{q}$ and divided by
  $\mu(X)^{1/q}$.
\item Consequently $N_q\to M$ as $q\to\infty$, and $q\mapsto N_q$ is
  nondecreasing.
\end{enumerate}
Clauses (1) and (2) are the Mathlib content; clause (3) is recorded here, rather
than left to each consumer, because it is the form actually used and it follows
from (1) and (2) in two lines: (1) gives
$\limsup_{q\to\infty}N_q\le M$, and for each $t<M$ clause (2) gives
$\liminf_{q\to\infty}N_q\ge t\lim_{q\to\infty}c^{1/q}=t$, where
$c=\mu\{|f|>t\}/\mu(X)\in(0,1]$, so $\liminf_{q\to\infty}N_q\ge M$ on letting
$t\uparrow M$; this covers $M=\infty$ as well, both sides then being $+\infty$.
The monotonicity in (3) is H\"older on a set of finite measure (Preamble)
applied to the normalised measure $\mu/\mu(X)$, which has total mass $1$: for
$q\le q'$ it gives $N_q\le N_{q'}$.
\end{proposition}

